\documentclass[12pt,a4paper]{article}  % Déclare la classe du document.

%  ------------------------------------------  Modules usuels  --------------------------------------
%  --------------------------------------------------------------------------------------------------------

\usepackage{amsmath} 	
\usepackage{amssymb} 
\usepackage{listings}         
\usepackage{numprint} 
\usepackage{pifont}

%  --------------------------------------------------------------------------------------------------------
% ----------------------  packages pour fontes et caractères spéciaux  ---------------------
%  --------------------------------------------------------------------------------------------------------

\usepackage{eurosym}   
\usepackage{amsfonts}   
\usepackage{amstext}       
\usepackage{amsfonts}	
\usepackage{mathrsfs}	 

\usepackage{pstricks,pst-all,pstricks-add,pst-tree,pst-node,pst-3dplot}

%------------------------------------------------------------------------------------------------------
%  --------------------------  Francisation et gestion des caractères sous mac --------------------------------------------------------------------------------------------------------

\usepackage[francais]{babel}   
\usepackage[utf8]{inputenc}
\usepackage[OT1,T1]{fontenc} 
%  --------------------------------------------------------------------------------------------------------
% ------------------------------------Packages  de  mise en page ---------------------------
%--------------------------------------------------------------------------------------------------------

\usepackage{fancyhdr}	% en-têtes et pieds de pages
\usepackage{amsfonts}

 \usepackage{multicol} % permet l'utilisation du multicolonnage dans le texte. La syntaxe est \begin{multicols}{nb de colonnes}    ......  \end{multicols}
 
\usepackage{calc}            % permet de faire un minimum de calcul sur les arguments des commandes et macros (utilisé dans la def de l'environnement preuve )
\usepackage{epic}            % extension et simplification de  l'environnement picture
\usepackage{graphicx}    % permet d'inclure des graphique externe dans un document
\usepackage{variations}   % permet d'obtenir de joli tableau de variations
\usepackage{enumerate}  %permet de faire des listes personnalisées
\usepackage{arydshln}     %permet d'avoir des lignes en pointillés
\usepackage{framed}  
\usepackage{multido} 


%  --------------------------------------  MACROS USUELLES  ----------------------------------------

\newcommand{\R}{\mathbb{R}} 
\newcommand{\N}{\mathbb{N}} 
\newcommand{\D}{\mathbb{D}} 
\newcommand{\Z}{\mathbb{Z}} 
\newcommand{\Q}{\mathbb{Q}} 
\newcommand{\C}{\mathbb{C}} 

%  -------------------------------------  MACROS SPECIFIQUES ------------------------------------


\newcommand{\raisedrule}[2][0em]{\leaders\hbox{\rule[#1]{1pt}{#2}}\hfill}


\newcounter{numexo}

\setcounter{numexo}{1}

\newcommand{\exercice}[1]{\vskip 0.6cm\noindent\rule[+2pt]{1cm}{0.5pt}\ {\large\bf \ Exercice \thenumexo {  #1 }}\raisedrule[+2pt]{0.5pt}\null\addtocounter{numexo}{1}\par\vskip 0,4cm}


\newcommand{\lignes}[1]{
\setlength{\unitlength}{1mm}
\vskip 0,4cm\noindent
\multido{\i=0+1}{#1}{%
     \noindent
     \begin{picture}(150,5)
     \linethickness{0,005mm}
     \put(-5,0){\dashline{1}(175,5)(5,5)}
     \end{picture}\vskip 0,001cm
     }
\vskip 0cm
}




% ----------------------------------------- FORMAT DE PAGE  -----------------------------------------

\usepackage[left = 1.3cm,
	 	      right = 1.3cm,
		      top = 1.7cm,
		      headsep = 8mm,
		      bottom = 1.7cm,
		      headheight = 14.2 pt
		      ]{geometry}

\setlength{\parindent}{0mm}


\everymath{\displaystyle} 
\renewcommand{\thesection}{\Roman{section}} 
\author{}
\title{}
\date{}
\pagestyle{fancy}

 
\lhead{1ère - spécialité maths} 


\rhead{\textcolor{red}{Correction - DS n$^{\circ}$6} - 26/04/21}		


\fancyfoot[C]{} 
\fancyfoot[L]{} 
\fancyfoot[R]{} 

\lstset{language=Python}  

\lstset{frame = single
framerule = 1pt
showspaces = false,
showstringspaces = false,
numbers=left, numberstyle=\tiny}

% ----------------------------------------- correction  -----------------------------------------
\usepackage[most,skins,breakable]{tcolorbox}
\usepackage{enumitem}

\newcommand{\corr}[1]{\fcolorbox{red}{yellow!30}{#1}}

\newcommand{\bareme}[1]{\tikz\node[circle,draw,inner sep=1pt,outer sep=1pt, color=red]{\textbf{#1}};}

\newtcolorbox{corrBox}{enhanced,boxrule=1pt,nobeforeafter,boxsep=0pt,left=\fboxsep,right=\fboxsep,colback=yellow!30,colframe=red!75!black}

\newlist{enumMet}{enumerate}{2}
\setlist[enumMet,1]{label=\arabic*), font=\bf\color{blue!75!black}}
\setlist[enumMet,2]{label=\alph*), font=\bf\color{blue!75!black}}


\begin{document}
\graphicspath{{images/}}

\begin{center}

\fbox{\LARGE{\rule{0.5cm}{0pt}\rule[-0.25cm]{0pt}{0.9cm} \textcolor{red}{Correction - DS n$^{\circ}$6 - Suites numériques}	\rule{0.5cm}{0pt}}}
\end{center}


\begin{center}\textbf{2h - Calculatrice autorisée - Barème indicatif}\end{center}

%\begin{center}\textbf{Les élèves avec un tier-temps ne traitent pas les questions avec  le symbole \ding{46}}\end{center}

\begin{center}
%On veillera à ENCADRER ses résultats et soigner sa copie.\\
Les résultats doivent être justifiés par des calculs (au moins 1 étape intermédiaire).
\end{center}
\vspace{.05cm} 

\vspace{-0.3cm}
\exercice{ - Reconnaître une suite arithmétique - ( 3 points )}

Reconnaître parmi les suites définies sur $\mathbb{N}$ ci-dessous celles qui sont arithmétiques et préciser alors leur premier terme et leur raison.
\begin{enumMet}
    \item $u_{n}=-2+3 n$ \quad\bareme{0,5}\\
    \begin{corrBox}
        On reconnaît une expression d'une suite arithmétique : $u_n=u_0+n\times r$ avec $u_0=-2$ et $r=3$
    \end{corrBox}
    \item $u_{n}=2 n^{2}+3$ \quad\bareme{0,5}\\
    \begin{corrBox}
        $u_0=3\quad\quad\quad u_1=5\quad\quad\quad u_2=11\quad\quad\quad\quad\quad\quad\quad\quad u_2-u_1=6\neq u_1-u_0=2$\\
        $u_n$ n'est donc pas une suite arithmétique.
    \end{corrBox}
    \item $u_{n}=\dfrac{n+5}{2}$ \quad\bareme{0,5}\\
    \begin{corrBox}
        On reconnaît une expression d'une suite arithmétique : $u_n=u_0+n\times r$ avec $u_0=\dfrac{5}{2}$ et $r=\dfrac{1}{2}$
    \end{corrBox}
    \item $u_{n}=1-\frac{1}{n+1}$ \quad\bareme{0,5}\\
    \begin{corrBox}
        $u_0=0\quad\quad\quad u_1=\dfrac{1}{2}\quad\quad\quad u_2=\dfrac{2}{3}\quad\quad\quad\quad\quad\quad\quad\quad u_2-u_1=\dfrac{1}{6}\neq u_1-u_0=\dfrac{1}{2}$\\
        $u_n$ n'est donc pas une suite arithmétique.
    \end{corrBox}
    \item $u_{0}=-1$ et $u_{n+1}=u_{n}+2 n$ \quad\bareme{0,5}\\
    \begin{corrBox}
        $u_{n+1}-u_n=u_n+2 n-u_n=2n$. La différence n'est pas constante. La suite $u_n$ n'est donc pas arithmétique. 
    \end{corrBox}
    \item $u_{0}=4$ et $u_{n+1}=1,5+u_{n}$ \quad\bareme{0,5}\\
    \begin{corrBox}
        $u_{n+1}-u_n=1,5+u_n-u_n=1,5=r$. La suite $u_n$ est une suite arithmétique de raison $r=1,5$ et de premier terme $u_0=4$. 
    \end{corrBox}
\end{enumMet}

\exercice{ - Reconnaître une suite géométrique - ( 2 points )}
Reconnaître parmi les suites définies sur $\mathbb{N}$ ci-dessous celles qui sont géométriques et préciser alors leur premier terme et leur raison.
\begin{enumMet}
    \item $u_{n}=4+n \times 4$\\
    \begin{corrBox}
        $u_0=4\quad\quad\quad u_1=\dfrac{17}{4}\quad\quad\quad u_2=\dfrac{9}{2}\quad\quad\quad\quad\quad\quad\quad\quad \dfrac{u_2}{u_1}=\dfrac{18}{17}\neq \dfrac{u_1}{u_0}=\dfrac{17}{8}$\\
        $u_n$ n'est donc pas une suite géométrique.\quad\bareme{0,5}
    \end{corrBox}
    \item $u_{n}=\frac{2^{n}}{3}$\\
    \begin{corrBox}
        $\dfrac{u_{n+1}}{u_n}=\dfrac{\dfrac{2^{n+1}}{3}}{\dfrac{2^n}{3}}=2$\\
        $u_n$ est une suite géométrique de raison $q=2$ et de premier terme $u_0=\dfrac{1}{3}$.\quad\bareme{0,5}
    \end{corrBox}
    \item $u_{n}=(\sqrt{2})^{n}$\\
    \begin{corrBox}
        $\dfrac{u_{n+1}}{u_n}=\dfrac{(\sqrt{2})^{n+1}}{(\sqrt{2})^{n}}=\sqrt{2}$\\
        $u_n$ est une suite géométrique de raison $q=\sqrt{2}$ et de premier terme $u_0=1$.\quad\bareme{0,5}
    \end{corrBox}
    \item $u_{1}=100$ et $u_{n+1}=\frac{u_{n}}{5}$\\
    \begin{corrBox}
        $u_{n+1}=\dfrac{1}{5}\times u_n$\\
        $u_n$ est une suite géométrique de raison $q=\dfrac{1}{5}$ et de premier terme $u_0=100\div \dfrac{1}{5}=500$.\quad\bareme{0,5}
    \end{corrBox}
\end{enumMet}

\exercice{ - Calculer des sommes - ( 4 points )}

\begin{enumMet}
    \item $S = 1+2+3+\ldots+1000$ \corr{$=\dfrac{1000\times1001}{2}=500500$}\quad\bareme{1}
    \item $T = 1+1,5+1,5^{2}+\ldots+1,5^{8}$ \corr{$=\dfrac{1-1,5^9}{1-1,5}=\dfrac{19171}{256}$}\quad\bareme{1}
    \item $Z = 501+502+503+\ldots+1000$ \corr{$=\dfrac{1000\times1001}{2}-\dfrac{500\times501}{2}=500500-125250=375250$}\quad\bareme{1}
    \item $R = 1-2+2^2+\ldots+(-2)^7$ \corr{$=\dfrac{1-(-2)^8}{1-(-2)}=-85$}\quad\bareme{1}
\end{enumMet}


\exercice{ - retrouver raison et premier terme - ( 2 points )}
\begin{enumMet}
    \item On a $u_{7}=4$ et $u_{4}=7$. Calculer $r$ puis $u_{35}$\\
    \begin{corrBox}
        $u_7-u_4=4-7 \iff u_0+7r-(u_0+4r)=-3 \iff 3r=-3 \iff r=-1$ \quad\bareme{0,5}\\
        $u_4=7\iff u_0+4\times(-1)=7\iff u_0-4=7 \iff u_0=11$\\
        $u_{35}=u_0+35r=11-35=-24$ \quad\bareme{0,5}
    \end{corrBox}
    \item On a $u_{9}=33$ et $u_{4}=13$. Calculer $r$ puis $u_{25}$.\\
    \begin{corrBox}
        $u_9-u_4=33-13 \iff u_0+9r-(u_0+4r)=20 \iff 5r=20 \iff r=4$ \quad\bareme{0,5}\\
        $u_4=13\iff u_0+4\times4=13\iff u_0+16=13 \iff u_0=-3$\\
        $u_25=u_0+25r=-3+25\times4=97$ \quad\bareme{0,5}
    \end{corrBox}
\end{enumMet}

\exercice{ - retrouver raison et premier terme - ( 2 points )}
\begin{enumMet}
    \item On a $u_{7}=11$ et $u_{10}=3$ 773. Calculer $q$ puis $u_{12}$.\\
    \begin{corrBox}
        $\dfrac{u_{10}}{u_7}=\dfrac{3773}{11}\iff \dfrac{q^{10}}{q^7}=343\iff q^3=343 \iff q=\sqrt[3]{343}\iff q=7$ \quad\bareme{0,5}\\
        $u_{12}=u_{10}\times q^2=3773\times 7^2 = 184\ 877$ \quad\bareme{0,5}
    \end{corrBox}
    \item On a $u_{5}=41$ et $u_{9}=25$ 625. Calculer $q$ puis $u_{10}$.\\
    \begin{corrBox}
        $\dfrac{u_{9}}{u_5}=\dfrac{25625}{41}\iff \dfrac{q^{9}}{q^5}=625\iff q^4=625 \iff q=\sqrt[4]{625}\iff q=5$ \quad\bareme{0,5}\\
        $u_{10}=u_{9}\times q=25625\times 5 = 128\ 125$ \quad\bareme{0,5}
    \end{corrBox}
\end{enumMet}


\exercice{ - sens de variation - ( 3 points )}
Déterminer le sens de variation des suites définies pour tout entier naturel $n$ par les formules explicites suivantes.
\begin{enumMet}
    \item $u_{n}=2-4 n$\\
    \begin{corrBox}
        $u_{n+1}-u_n=-4<0$. La suite est donc décroissante \quad\bareme{0,5}
    \end{corrBox}
    \item $v_{n}=2 n^{2}+3$\\
    \begin{corrBox}
        $v_{n+1}-v_n=2n^2+4n+2-2n^2=4n+2>0$ pour $n \in \Z$. La suite est donc croissante. \quad\bareme{0,5}
    \end{corrBox}
    \item $w_{n}=n^{2}+2 n$\\
    \begin{corrBox}
        $w_{n+1}-w_n=n^2+2n+1+2n+2-n^2-2n=2n+3>0$ pour $n \in \Z$. La suite est donc croissante. \quad\bareme{0,5}
    \end{corrBox}
    \item $t_{n}=2^{n}+3^{n}$\\
    \begin{corrBox}
        $t_{n+1}-t_n=2^{n+1}+3^{n+1}-2^n-3^n=2^n(2-1)+3^n(3-1)=2^n+2\times3^n>0$ pour $n \in \Z$. La suite est donc croissante. \quad\bareme{0,5}
    \end{corrBox}
    \item $x_{n}=n^{2}-8 n+2$\\
    \begin{corrBox}
        $x_{n+1}-x_n=n^2+2n+1-8n-8+2-n^2+8n-2=2n-7>0$ pour $n>3,5$. La suite est donc décroissante sur $[0;3]$ puis croissante sur $[4;+\infty]$. \quad\bareme{1}
    \end{corrBox}
\end{enumMet}

\newpage

\exercice{ - sens de variation d'une suite arithmétique - ( 2 points )}
Pour les suites arithmétiques suivantes dont on donne $\mathrm{l} \mathrm{}1^{\textrm{er}}$ terme et la raison, déterminer le sens de variation.
\begin{enumMet}
    \item $u_{0}=-2$ et $r=0,6$\\
    \begin{corrBox}
        $r>0$, la suite est donc croissante. \quad\bareme{0,5}
    \end{corrBox}
    \item $v_{0}=1$ et $r=\frac{2}{3}$.\\
    \begin{corrBox}
        $r>0$, la suite est donc croissante. \quad\bareme{0,5}
    \end{corrBox}
    \item $w_{0}=5$ et $r=1-\sqrt{2}$.\\
    \begin{corrBox}
        $r<0$, la suite est donc décroissante. \quad\bareme{0,5}
    \end{corrBox}
    \item $t_{0}=-10$ et $r=10^{-2}$.\\
    \begin{corrBox}
        $r>0$, la suite est donc croissante. \quad\bareme{0,5}
    \end{corrBox}
\end{enumMet}

\exercice{ - sens de variation d'une suite géométrique - ( 2 points )}
Pour les suites géométriques suivantes dont on donne le $1^{\text {er }}$ terme et la raison, déterminer le sens de variation
\begin{enumMet}
    \item $u_{0}=3$ et $q=2$.\\
    \begin{corrBox}
        $u_0>0$ et $q>1$, la suite est donc croissante. \quad\bareme{0,5}
    \end{corrBox}
    \item $v_{0}=-1$ et $q=\frac{4}{5}$.\\
    \begin{corrBox}
        $v_0<0$ et $q<1$, la suite est donc croissante. \quad\bareme{0,5}
    \end{corrBox}
    \item $w_{0}=\frac{-2}{3}$ et $q=\frac{8}{3}$.\\
    \begin{corrBox}
        $w_0<0$ et $q>1$, la suite est donc décroissante. \quad\bareme{0,5}
    \end{corrBox}
    \item $t_{0}=0,5$ et $q=10^{-1}$.\\
    \begin{corrBox}
        $t_0>0$ et $q<1$, la suite est donc décroissante. \quad\bareme{0,5}
    \end{corrBox}
\end{enumMet}

\end{document}

